% FIG ch06-04 structure errors (Chinese)%
% Shared styles for chapter 6 workflow TikZ figures.%
% Every figure inputs this file, so commands are defined with \providecommand.%
\providecolor{BookTeal}{HTML}{176C68}%
\providecolor{BookInk}{HTML}{23363E}%
\providecolor{BookRule}{HTML}{C8D6D6}%
\providecolor{BookPale}{HTML}{F1F6F5}%
\definecolor{WetPale}{HTML}{FBF1EA}%
\definecolor{Wet}{HTML}{B5653A}%
\definecolor{Seq}{HTML}{8A6D1F}%
\definecolor{DryPale}{HTML}{EEF4F7}%
\definecolor{Dry}{HTML}{2F6F8F}%
\definecolor{ChSixFeedback}{HTML}{7A8A90}%
\definecolor{ChSixBad}{HTML}{C0392B}%
\usetikzlibrary{arrows.meta,positioning,shapes.geometric,shapes.misc,fit,backgrounds,calc}%
\ifdefined\ChSixLayersReady\else%
  \pgfdeclarelayer{chsixlane}%
  \pgfsetlayers{chsixlane,background,main}%
  \def\ChSixLayersReady{}%
\fi%

\tikzset{%
  chsix/.style={font=\sffamily\scriptsize, text=BookInk, align=center},%
  flowstep/.style={%
    chsix, rounded corners=2pt, draw=Wet, line width=0.45pt, fill=white,%
    text width=19mm, minimum height=7.4mm, inner sep=1.3pt,%
  },%
  flowstepdry/.style={flowstep, draw=Dry},%
  flowstepseq/.style={flowstep, draw=Seq},%
  artifact/.style={%
    chsix, trapezium, trapezium left angle=70, trapezium right angle=110,%
    trapezium stretches=true,%
    draw=BookInk, line width=0.45pt, fill=white,%
    minimum height=6.2mm, inner sep=1.2pt,%
  },%
  gate/.style={%
    chsix, diamond, aspect=1.7, draw=BookInk, line width=0.45pt, fill=white,%
    text width=12mm, inner sep=0.6pt, minimum height=8mm,%
  },%
  opt/.style={flowstep, dashed, draw=BookInk!70},%
  optdry/.style={flowstepdry, dashed, draw=BookInk!70},%
  stack/.style={flowstep},%
  link/.style={-{Stealth[length=1.6mm]}, semithick, draw=BookInk},%
  linkthin/.style={-{Stealth[length=1.2mm]}, thin, draw=BookInk!80},%
  back/.style={-{Stealth[length=1.5mm]}, dashed, draw=ChSixFeedback},%
  conn/.style={%
    circle, draw=BookInk, line width=0.5pt, minimum size=4.6mm,%
    font=\sffamily\scriptsize\bfseries, inner sep=0pt, text=BookInk,%
  },%
  lanetitle/.style={font=\sffamily\bfseries\small, text=BookInk, align=left},%
  passlabel/.style={font=\sffamily\scriptsize, text=BookInk, fill=white, inner sep=0.7pt, align=center},%
  mergelabel/.style={font=\sffamily\scriptsize, text=BookInk, fill=white, inner sep=0.8pt, align=center},%
  taglab/.style={font=\sffamily\scriptsize, text=ChSixFeedback, fill=white, inner sep=0.8pt, align=center},%
  gramcap/.style={font=\sffamily\scriptsize, text=BookInk, align=left, anchor=west},%
}%

% Offset card behind a two-sample step. Text stays on the front node only.%
\providecommand{\ChSixCardBack}[1]{%
  \begin{pgfonlayer}{background}%
    \draw[Wet, line width=0.45pt, rounded corners=2pt, fill=white]%
      ([shift={(-0.55mm,0.55mm)}]#1.south west)%
      rectangle ([shift={(-0.55mm,0.55mm)}]#1.north east);%
  \end{pgfonlayer}%
}%
% #1 options, #2 name, #3 position (with parentheses), #4 text%
\providecommand{\ChSixCard}[4]{%
  \node[stack, #1] (#2) at #3 {#4};%
  \ChSixCardBack{#2}%
}%

\providecommand{\ChSixMaybeLabel}[1]{%
  \ifdefined\nolabels\else#1\fi%
}%

\providecommand{\ChSixIcon}[2]{%
  \ifdefined\ChSixShowIcons%
    \IfFileExists{#1}{%
      \node[anchor=south, inner sep=0] at ([yshift=0.7mm]#2.north)%
        {\includegraphics[height=5.4mm]{#1}};%
    }{}%
  \fi%
}%
%
\begin{tikzpicture}[x=1mm, y=1mm, font=\sffamily\scriptsize, text=BookInk]%
  \tikzset{%
    onestep/.style={flowstep, minimum height=6.5mm},%
    onedry/.style={flowstepdry, minimum height=6.5mm},%
    oneseq/.style={flowstepseq, minimum height=6.5mm},%
    twostep/.style={flowstep, minimum height=8.4mm},%
    twoseq/.style={flowstepseq, minimum height=8.4mm},%
    prod/.style={artifact, minimum height=6.5mm},%
  }%
  % (a)%
  \node[anchor=north west, font=\sffamily\bfseries] (aH) at (0,0)%
    {(a)\ 可视化被当成分析的上游};%
  \node[onedry, anchor=north west] (a1b)%
    at ([xshift=11mm, yshift=-1.5mm]aH.south west) {k 近邻图};%
  \node[onedry, anchor=west] (a2b) at ([xshift=6mm]a1b.east) {UMAP};%
  \node[onedry, anchor=west] (a3b) at ([xshift=6mm]a2b.east) {聚类};%
  \node[onedry, anchor=west] (a4b) at ([xshift=6mm]a3b.east) {注释};%
  \draw[link] (a1b) -- (a2b) -- (a3b) -- (a4b);%
  \node[anchor=east, font=\sffamily\bfseries, text=ChSixBad]%
    at ([xshift=-4mm]a1b.west) {\bookcross};%

  \node[onedry, anchor=north west] (a1g) at ([yshift=-2.6mm]a1b.south west) {k 近邻图};%
  \node[onedry, anchor=west] (a2g) at ([xshift=6mm]a1g.east) {聚类};%
  \node[onedry, anchor=west] (a3g) at ([xshift=6mm]a2g.east) {注释};%
  \node[onedry, anchor=north] (a4g) at ([yshift=-3.2mm]a1g.south) {UMAP};%
  \node[prod, anchor=west] (a5g) at ([xshift=6mm]a4g.east) {仅用于展示};%
  \draw[link] (a1g) -- (a2g) -- (a3g);%
  \draw[link] (a1g.south) -- (a4g.north);%
  \draw[link] (a4g) -- (a5g);%
  \coordinate (aGmid) at ($(a1g.north)!0.5!(a4g.south)$);%
  \coordinate (aMark) at (a1g.west |- aGmid);%
  \node[anchor=east, font=\sffamily\bfseries, text=BookTeal]%
    at ([xshift=-4mm]aMark) {\bookcheck};%
  \coordinate (aNoteAt) at (a4g.south -| 0,0);%
  \node[anchor=north west, text width=120mm, text=BookInk!75] (aNote)%
    at ([yshift=-1.5mm]aNoteAt)%
    {把 UMAP 改为从 k 近邻图分出的平行分支，不连向聚类和注释};%

  % (b)%
  \node[anchor=north west, font=\sffamily\bfseries] (bH)%
    at ([yshift=-1.8mm]aNote.south west) {(b)\ 步骤与产物混在一个框};%
  \node[onedry, anchor=north west, text width=42mm] (b1)%
    at ([xshift=11mm, yshift=-1.5mm]bH.south west) {FASTQ 比对矩阵};%
  \node[anchor=east, font=\sffamily\bfseries, text=ChSixBad]%
    at ([xshift=-4mm]b1.west) {\bookcross};%

  \node[prod, anchor=north west, minimum width=16mm] (b2a)%
    at ([yshift=-2.8mm]b1.south west) {FASTQ};%
  \node[onedry, anchor=west, text width=32mm] (b2b)%
    at ([xshift=6.5mm]b2a.east) {比对与 UMI 计数};%
  \node[prod, anchor=west, text width=30mm] (b2c)%
    at ([xshift=6.5mm]b2b.east) {原始计数矩阵};%
  \draw[link] (b2a) -- (b2b);%
  \draw[link] (b2b) -- (b2c);%
  \node[anchor=east, font=\sffamily\bfseries, text=BookTeal]%
    at ([xshift=-4mm]b2a.west) {\bookcheck};%
  \coordinate (bNoteAt) at (b2b.south -| 0,0);%
  \node[anchor=north west, text width=120mm, text=BookInk!75] (bNote)%
    at ([yshift=-1.5mm]bNoteAt)%
    {步骤框与产物平行四边形分开，交接才可见};%

  % (c)%
  \node[anchor=north west, font=\sffamily\bfseries] (cH)%
    at ([yshift=-1.8mm]bNote.south west) {(c)\ 质控没有出口};%
  \node[onestep, anchor=north west, text width=26mm] (c1b)%
    at ([xshift=11mm, yshift=-1.5mm]cH.south west) {解离};%
  \node[onestep, anchor=west, text width=30mm] (c2b)%
    at ([xshift=6mm]c1b.east) {活率检查};%
  \node[onestep, anchor=west, text width=20mm] (c3b)%
    at ([xshift=6mm]c2b.east) {上样};%
  \draw[link] (c1b) -- (c2b) -- (c3b);%
  \node[anchor=east, font=\sffamily\bfseries, text=ChSixBad]%
    at ([xshift=-4mm]c1b.west) {\bookcross};%

  \node[onestep, anchor=north west, text width=26mm] (c1g)%
    at ([yshift=-3.6mm]c1b.south west) {解离};%
  \node[gate, anchor=west, text width=15mm, minimum height=9mm, aspect=1.6] (c2g)%
    at ([xshift=6.5mm]c1g.east) {活率};%
  \node[onestep, anchor=west, text width=20mm] (c3g)%
    at ([xshift=13mm]c2g.east) {上样};%
  \draw[link] (c1g) -- (c2g);%
  \draw[link] (c2g) -- (c3g);%
  \node[passlabel, anchor=south] (cPass)%
    at ([yshift=1.8mm]$(c2g.east)!0.5!(c3g.west)$) {通过};%
  \coordinate (cBase) at ($(c1g.south)!0.5!(c2g.south)+(0,-3.8mm)$);%
  \draw[back, rounded corners=1.6mm] (c2g.south) -- (c2g.south |- cBase)%
    -- (c1g.south |- cBase) -- (c1g.south);%
  \node[taglab, anchor=north] (cFailLab) at ([yshift=-1.2mm]cBase) {不通过};%
  \coordinate (cMid) at ($(c2g.north)!0.5!(cFailLab.south)$);%
  \coordinate (cMark) at (c1g.west |- cMid);%
  \node[anchor=east, font=\sffamily\bfseries, text=BookTeal]%
    at ([xshift=-4mm]cMark) {\bookcheck};%
  \coordinate (cNoteAt) at (cFailLab.south -| 0,0);%
  \node[anchor=north west, text width=120mm, text=BookInk!75] (cNote)%
    at ([yshift=-1.5mm]cNoteAt)%
    {菱形 + 不通过时返回上一步};%

  % (d)%
  \node[anchor=north west, font=\sffamily\bfseries] (dH)%
    at ([yshift=-1.8mm]cNote.south west) {(d)\ 并行被画成串联};%
  \node[twostep, anchor=north west, text width=28mm] (d1b)%
    at ([xshift=11mm, yshift=-1.5mm]dH.south west) {样本 1\\解离};%
  \node[twostep, anchor=west, text width=28mm] (d2b)%
    at ([xshift=6mm]d1b.east) {样本 2\\解离};%
  \node[twoseq, anchor=west, text width=24mm] (d3b)%
    at ([xshift=6mm]d2b.east) {测序};%
  \draw[link] (d1b) -- (d2b) -- (d3b);%
  \node[anchor=east, font=\sffamily\bfseries, text=ChSixBad]%
    at ([xshift=-4mm]d1b.west) {\bookcross};%

  \node[twostep, anchor=north west, text width=40mm] (d1g)%
    at ([yshift=-2.6mm]d1b.south west) {样本 1：\\解离 $\rightarrow$ 文库};%
  \node[twostep, anchor=north west, text width=40mm] (d2g)%
    at ([yshift=-1.8mm]d1g.south west) {样本 2：\\解离 $\rightarrow$ 文库};%
  \coordinate (dJoin) at ([xshift=7mm]$(d1g.east)!0.5!(d2g.east)$);%
  \node[twoseq, anchor=west, text width=26mm] (d3g)%
    at ([xshift=5.5mm]dJoin) {混合与测序};%
  \draw[link] (d1g.east) -- (d1g.east -| dJoin);%
  \draw[link] (d2g.east) -- (d2g.east -| dJoin);%
  \draw[BookInk, semithick] (d1g.east -| dJoin) -- (d2g.east -| dJoin);%
  \fill[BookInk] (dJoin) circle (0.7pt);%
  \draw[link] (dJoin) -- (d3g.west);%
  \coordinate (dMid) at ($(d1g.north)!0.5!(d2g.south)$);%
  \coordinate (dMark) at (d1g.west |- dMid);%
  \node[anchor=east, font=\sffamily\bfseries, text=BookTeal]%
    at ([xshift=-4mm]dMark) {\bookcheck};%
  \coordinate (dNoteAt) at (d2g.south -| 0,0);%
  \node[anchor=north west, text width=120mm, text=BookInk!75]%
    at ([yshift=-1.5mm]dNoteAt)%
    {两样本并列处理，在合并点汇合后再测序};%
\end{tikzpicture}%
